\par\addvspace{12pt}\Needspace{9\baselineskip}

\section{函数的可积性*}

% ITEM: calculus-example-056
\solutionitem{calculus-example-056}{例5.56 *}
对\(f(x)=x^2\)，求\([0,1]\)上等分分割的上下和之差。再考察只在\(x=1/2\)处取\(1\)、其余位置为零的函数：哪些小段有非零振幅？说明这些小段的总贡献为何趋零。
\solutionanswer
\(x^2\)的上和减下和为\(n^{-3}\sum_{i=1}^n[i^2-(i-1)^2]=1/n\)。单点非零函数只有接触\(1/2\)的一段或两段振幅为\(1\)，其余振幅为零；因此\(U-L\le2\|P\|\to0\)，可积且积分零。


% ITEM: calculus-exercise-34-01
\solutionitem{calculus-exercise-34-01}{习题5.35（1） *}
设\(|f(u)-f(v)|\le C\sqrt{|u-v|}\)。给出上下和之差关于\(\|P\|\)的界，并证明可积性。
\solutionanswer
每段振幅\(\omega_i\le C\sqrt{\Delta x_i}\le C\sqrt{\|P\|}\)，于是\(U-L\le C(b-a)\sqrt{\|P\|}\to0\)。条件也使\(f\)有界，Darboux判据给出可积性。


% ITEM: calculus-exercise-34-02
\solutionitem{calculus-exercise-34-02}{习题5.35（2） *}
用上题的方法处理\(f(x)=\sqrt{x}\)、\(0\le x\le1\)。为什么不能直接在整个区间上使用一个有限的导数上界？
\solutionanswer
\(|\sqrt u-\sqrt v|\le\sqrt{|u-v|}\)，故\(U-L\le\sqrt{\|P\|}\to0\)。\(f'(x)=1/(2\sqrt x)\)在零点附近无界，不能使用整个区间统一的有限导数界；函数值的平方根变化界却仍有效。


% ITEM: calculus-exercise-34-03
\solutionitem{calculus-exercise-34-03}{习题5.35（3） *}
对\(f(x)=|x-1/2|\)、\(0\le x\le1\)，估计\(n\)等分时的上下和之差，并求使任意两组取样和相差小于\(10^{-3}\)的一个充分条件。
\solutionanswer
绝对值三角不等式给出\(|f(u)-f(v)|\le|u-v|\)。在\(n\)等分下\(U-L\le1/n\)，所以任意两组取样和之差不超过\(1/n\)。取整数\(n\ge1001\)即可严格小于\(10^{-3}\)。


% ITEM: calculus-example-057
\solutionitem{calculus-example-057}{例5.57 *}
本题的阶跃函数为\(x<1/2\)时\(0\)、\(x\ge1/2\)时\(1\)，定义于\([0,1]\)。本题中\(g(0)=0\)，且当\(1/(n+1)<x\le1/n\)、\(n\ge1\)时\(g(x)=1/n\)。

对上一节的阶跃函数和无穷阶梯函数\(g\)，分别解释上述估计如何证明可积性。是否需要先数清全部间断点？
\solutionanswer
阶跃函数单调不减，总增量为\(1\)，故\(U-L\le\|P\|\)。无穷阶梯函数\(g\)也单调不减，且\(g(1)-g(0)=1\)，同样有\(U-L\le\|P\|\)。两者都由Darboux判据可积，不必逐个枚举跳点；控制全部小段增量的总和已经足够。


% ITEM: calculus-exercise-35-01
\solutionitem{calculus-exercise-35-01}{习题5.36（1） *}
对\(f(x)=\sin(1/x)\)、\(x>0\)，零点补零，具体说明怎样先选\([0,\ell]\)，再在\([\ell,1]\)上细分，使上下和之差小于给定误差。
\solutionanswer
给定\(\varepsilon>0\)，先取\(0<\ell<\min(1,\varepsilon/4)\)，使第一段贡献\(2\ell<\varepsilon/2\)。在\([\ell,1]\)，导数绝对值不超过\(1/\ell^2\)，故将该段细分至最大段长\(\delta<\varepsilon\ell^2/[2(1-\ell)]\)（若剩余段非空），其上下和之差小于\(\varepsilon/2\)。总差小于\(\varepsilon\)。


% ITEM: calculus-exercise-35-02
\solutionitem{calculus-exercise-35-02}{习题5.36（2） *}
某函数除有限个点外等于\(x^2\)。证明它在\([0,1]\)上可积并求积分。若改为所有有理点处任意改变取值，原结论是否还必然成立？
\solutionanswer
与\(x^2\)之差只在有限\(N\)点非零，设该差绝对值不超过\(B\)，则任意取样和的差不超过\(2NB\|P\|\to0\)。因此可积，积分仍为\(1/3\)。若在全部有理点将值改为\(x^2+1\)，保持无理点值为\(x^2\)，每段上有至少\(1\)的振幅，便不可积；无限稠密的改变不能照搬有限点的论证。


% ITEM: calculus-exercise-36-01
\solutionitem{calculus-exercise-36-01}{习题5.37（1） *}
若\(f,g\)可积，证明\(|f-g|\)、\((f+g)^2\)可积。说明积分\(\int fg\)与\((\int f)(\int g)\)一般没有相等关系。
\solutionanswer
线性组合\(f-g,f+g\)可积，再由绝对值与乘积可积性，\(|f-g|\)、\((f+g)^2\)可积。在\([0,1]\)取\(f=g=x\)，则\(\int fg=1/3\)，而\((\int f)(\int g)=1/4\)，一般不相等。


% ITEM: calculus-exercise-36-02
\solutionitem{calculus-exercise-36-02}{习题5.37（2） *}
构造不可积函数\(h\)，使\(|h|\)和\(h^2\)都可积。可考虑在有理点、无理点取不同符号。
\solutionanswer
在有理点取\(h=1\)，无理点取\(h=-1\)。任意非零长度小段振幅均为\(2\)，所以不可积；但\(|h|=h^2=1\)是常值函数，都可积。


% ITEM: calculus-exercise-36-03
\solutionitem{calculus-exercise-36-03}{习题5.37（3） *}
本题\(D(x)=1\)（\(x\in\mathbb Q\)），\(D(x)=0\)（\(x\notin\mathbb Q\)），定义于\([0,1]\)。
对上一节的\(D\)，直接求任意分割的上和与下和，用判据再次说明它不可积。
\solutionanswer
每个小段均含有理数和无理数，所以\(M_i=1,m_i=0\)。任意分割均有\(U(D;P)=1,L(D;P)=0\)，差恒为\(1\)，不满足Darboux判据。
\par
